\documentclass[aspectratio=169]{beamer}
\usetheme{Madrid}
\usecolortheme{whale}
\usepackage[utf8]{inputenc}
\usepackage[spanish]{babel}
\usepackage{amsmath,amssymb}
\usepackage{graphicx}
\usepackage{booktabs}
\usepackage{listings}
\usepackage{xcolor}

\definecolor{codegreen}{rgb}{0,0.6,0}
\definecolor{codegray}{rgb}{0.5,0.5,0.5}
\definecolor{codepurple}{rgb}{0.58,0,0.82}
\definecolor{backcolour}{rgb}{0.95,0.95,0.92}

\lstdefinestyle{pythonstyle}{
    backgroundcolor=\color{backcolour},
    commentstyle=\color{codegreen},
    keywordstyle=\color{blue},
    numberstyle=\tiny\color{codegray},
    stringstyle=\color{codepurple},
    basicstyle=\ttfamily\footnotesize,
    breakatwhitespace=false,
    breaklines=true,
    captionpos=b,
    keepspaces=true,
    showspaces=false,
    showstringspaces=false,
    showtabs=false,
    tabsize=2
}

\lstset{style=pythonstyle}

\title{Diagnóstico del Modelo}
\subtitle{Verificación de Supuestos}
\author{Métodos de Análisis de Datos 1}
\institute{Universidad Nacional del Sur}
\date{Unidad 3 - Clase 16}

\begin{document}

\frame{\titlepage}

\begin{frame}{Contenidos}
\tableofcontents
\end{frame}

\section{Introducción}

\begin{frame}{¿Por qué diagnosticar?}
Ajustar un modelo es fácil, validarlo es crucial.

\vspace{0.5cm}
\textbf{Supuestos (L.I.H.N.):}
\begin{enumerate}
    \item \textbf{L}inealidad: $E(Y|X) = \beta_0 + \beta_1 X$
    \item \textbf{I}ndependencia: observaciones independientes
    \item \textbf{H}omocedasticidad: varianza constante
    \item \textbf{N}ormalidad: $\epsilon \sim N(0, \sigma^2)$
\end{enumerate}

\vspace{0.5cm}
\textbf{Si se violan:}
\begin{itemize}
    \item Predicciones sesgadas
    \item Errores estándar incorrectos
    \item Tests e intervalos inválidos
\end{itemize}
\end{frame}

\section{Gráficos de Residuos}

\begin{frame}{Los residuos}
\begin{block}{Residuo ordinario}
\[
e_i = Y_i - \hat{Y}_i
\]
\end{block}

\begin{block}{Residuo estandarizado}
\[
r_i = \frac{e_i}{s\sqrt{1 - h_i}}
\]
\end{block}

donde $h_i$ es el leverage y $s = \sqrt{\text{SCR}/(n-2)}$

\vspace{0.3cm}
Los residuos estandarizados tienen media 0 y varianza $\approx 1$
\end{frame}

\begin{frame}{Gráfico 1: Residuos vs Ajustados}
\textbf{Detecta:} No linealidad y heterocedasticidad

\vspace{0.3cm}
\textbf{Lo ideal:} Patrón aleatorio sin estructura

\vspace{0.3cm}
\textbf{Problemas:}
\begin{itemize}
    \item Patrón curvo $\Rightarrow$ relación no lineal
    \item Forma de embudo $\Rightarrow$ heterocedasticidad
\end{itemize}

\vspace{0.5cm}
\textbf{Soluciones:}
\begin{itemize}
    \item Transformaciones (log, raíz cuadrada)
    \item Regresión polinomial
    \item Errores estándar robustos
\end{itemize}
\end{frame}

\begin{frame}{Gráfico 2: Q-Q Plot}
\textbf{Detecta:} Violación de normalidad

\vspace{0.3cm}
\textbf{Lo ideal:} Puntos sobre la diagonal

\vspace{0.3cm}
\textbf{Desviaciones:}
\begin{itemize}
    \item Colas pesadas: puntos se alejan en extremos
    \item Asimetría: curvatura sistemática
\end{itemize}

\vspace{0.5cm}
\textbf{Nota:} La normalidad es menos crítica con $n$ grande (Teorema Central del Límite)
\end{frame}

\begin{frame}{Gráfico 3: Scale-Location}
\textbf{Detecta:} Heterocedasticidad de forma más clara

\vspace{0.3cm}
Grafica $\sqrt{|r_i|}$ vs valores ajustados

\vspace{0.3cm}
\textbf{Lo ideal:} Banda horizontal con dispersión uniforme

\vspace{0.3cm}
\textbf{Problema:} Tendencia ascendente/descendente $\Rightarrow$ varianza no constante
\end{frame}

\begin{frame}{Gráfico 4: Residuos vs Leverage}
\textbf{Detecta:} Observaciones influyentes

\vspace{0.3cm}
\textbf{Conceptos:}
\begin{itemize}
    \item \textbf{Leverage} $h_i$: qué tan lejos está $X_i$ de $\bar{X}$
    \[
    h_i = \frac{1}{n} + \frac{(X_i - \bar{X})^2}{S_{XX}}
    \]
    \item \textbf{Distancia de Cook} $D_i$: influencia total
    \[
    D_i = \frac{r_i^2}{p+1} \cdot \frac{h_i}{1-h_i}
    \]
\end{itemize}

\vspace{0.3cm}
\textbf{Regla:} $D_i > 1$ $\Rightarrow$ observación muy influyente
\end{frame}

\section{Tests Formales}

\begin{frame}{Test de normalidad}
\textbf{Shapiro-Wilk:}
\[
H_0: \text{residuos normales} \quad \text{vs} \quad H_1: \text{no normales}
\]

\vspace{0.5cm}
\textbf{Nota:} Con $n$ grande, muy sensible. Siempre complementar con Q-Q plot.
\end{frame}

\begin{frame}{Test de homocedasticidad}
\textbf{Breusch-Pagan:}
\[
H_0: \text{Var}(\epsilon) = \sigma^2 \quad \text{vs} \quad H_1: \text{varianza no constante}
\]

\vspace{0.5cm}
\textbf{Test de White:} Versión más robusta
\end{frame}

\begin{frame}{Test de independencia}
\textbf{Durbin-Watson:} Para series temporales
\[
DW = \frac{\sum_{i=2}^{n}(e_i - e_{i-1})^2}{\sum_{i=1}^{n} e_i^2}
\]

\begin{itemize}
    \item $DW \approx 2$: no autocorrelación
    \item $DW < 2$: autocorrelación positiva
    \item $DW > 2$: autocorrelación negativa
\end{itemize}
\end{frame}

\section{Implementación}

\begin{frame}[fragile]{Python: Gráficos}
\begin{lstlisting}[language=Python]
import matplotlib.pyplot as plt
import scipy.stats as stats

fig, axes = plt.subplots(2, 2, figsize=(12, 10))

# 1. Residuos vs Ajustados
axes[0,0].scatter(modelo.fittedvalues, modelo.resid)
axes[0,0].axhline(0, color='red', linestyle='--')
axes[0,0].set_title('Residuos vs Ajustados')

# 2. Q-Q Plot
stats.probplot(modelo.resid, dist="norm", 
               plot=axes[0,1])

# 3. Scale-Location
axes[1,0].scatter(modelo.fittedvalues, 
                  np.sqrt(np.abs(modelo.resid_pearson)))

# 4. Residuos vs Leverage
influence = modelo.get_influence()
axes[1,1].scatter(influence.hat_matrix_diag, 
                  modelo.resid_pearson)
plt.show()
\end{lstlisting}
\end{frame}

\begin{frame}[fragile]{Python: Tests}
\begin{lstlisting}[language=Python]
from scipy import stats
from statsmodels.stats.diagnostic import het_breuschpagan
from statsmodels.stats.stattools import durbin_watson

# Shapiro-Wilk
stat, p = stats.shapiro(modelo.resid)
print(f"Shapiro-Wilk p-valor: {p:.4f}")

# Breusch-Pagan
bp = het_breuschpagan(modelo.resid, modelo.model.exog)
print(f"Breusch-Pagan p-valor: {bp[1]:.4f}")

# Durbin-Watson
dw = durbin_watson(modelo.resid)
print(f"Durbin-Watson: {dw:.4f}")

# Cook's distance
cooks = influence.cooks_distance[0]
print(f"Obs influyentes: {np.where(cooks > 1)[0]}")
\end{lstlisting}
\end{frame}

\begin{frame}[fragile]{R: Todo en uno}
\begin{lstlisting}[language=R]
# Graficos automaticos
par(mfrow = c(2, 2))
plot(modelo)
par(mfrow = c(1, 1))

# Tests
shapiro.test(modelo$residuals)

library(lmtest)
bptest(modelo)
dwtest(modelo)

# Cook's distance
cooks <- cooks.distance(modelo)
plot(cooks, type="h")
abline(h=1, col="red", lty=2)
\end{lstlisting}
\end{frame}

\section{Resumen}

\begin{frame}{Resumen}
\textbf{4 gráficos clave:}
\begin{enumerate}
    \item Residuos vs Ajustados (linealidad, homocedasticidad)
    \item Q-Q Plot (normalidad)
    \item Scale-Location (homocedasticidad)
    \item Residuos vs Leverage (influencia)
\end{enumerate}

\vspace{0.3cm}
\textbf{Tests:} Shapiro-Wilk, Breusch-Pagan, Durbin-Watson

\vspace{0.3cm}
\textbf{Medidas:} Leverage, Distancia de Cook

\vspace{0.5cm}
\textbf{Próxima clase:} Regresión lineal múltiple
\end{frame}

\end{document}
